\section{Ejecuciones piloto}
\label{sec:piloto}

Se realizaron ejecuciones piloto con la configuración base con el objeto de
detectar comportamientos anómalos antes de las corridas de producción.

\subsection{Utilización observada}

\begin{table}[H]
\centering
\begin{tabular}{lccc}
\toprule
\textbf{Recurso} & \textbf{Servidores} & \textbf{$\rho$ analítico} & \textbf{$\rho$ observado} \\
\midrule
Triage       & 1 & 0,88 & 0,958 \\
Consultorios & 4 & 0,84 & 0,939 \\
\bottomrule
\end{tabular}
\caption{Utilización analítica y observada en las ejecuciones piloto.}
\label{tab:piloto}
\end{table}

\subsection{Hallazgo principal}

Las ejecuciones piloto revelaron que el sistema no alcanza régimen
estacionario dentro del horizonte considerado. Al cierre de la corrida se
contabilizaron \textbf{55 pacientes en espera de consultorio} frente a
\textbf{5 en espera de triage}, sobre un total aproximado de 72 pacientes
arribados.

Dos conclusiones se desprenden de esta observación:

\begin{enumerate}
  \item \textbf{La restricción dominante son los consultorios médicos}, no el
  puesto de triage. La cola se acumula aguas abajo de la clasificación.
  \item \textbf{La magnitud de la acumulación excede lo esperable} para una
  utilización de 0,84. Un sistema M/M/4 con esa carga presenta una longitud
  media de cola del orden de tres a cinco pacientes, no de cincuenta y cinco.
\end{enumerate}

\begin{pendiente}
El segundo punto indica que el modelo todavía contiene un desajuste no
resuelto. La hipótesis principal, ya señalada en la sección~1.4, es que la
exclusión de traumatología, radiografías y medicación concentra en cuatro
consultorios una demanda que el sistema real reparte entre más sectores. Esta
sección debe completarse una vez resuelto el punto, incorporando los valores
definitivos de la tabla~\ref{tab:piloto}.
\end{pendiente}

\subsection{Sesgo de truncamiento}

Las ejecuciones piloto pusieron de manifiesto un problema metodológico
relevante. Las métricas de espera se registraban en el momento del egreso del
paciente, de modo que los pacientes que permanecían en cola al finalizar la
corrida no quedaban contabilizados.

En un sistema congestionado, los pacientes que alcanzan a egresar son
precisamente aquellos que esperaron menos, lo que introduce un sesgo
sistemático hacia valores bajos. El sesgo es además de magnitud distinta
entre escenarios, lo que invalida la comparación: cuanto más congestionado el
escenario, mayor la proporción de pacientes truncados y menor el promedio
observado.

La corrección aplicada consiste en registrar los tiempos de espera en el
instante en que comienza la atención y no en el egreso, de modo que todo
paciente que llegó a ser atendido quede contabilizado con independencia de si
completó su recorrido dentro del horizonte. Adicionalmente se registra el
número de observaciones de cada estadística como control de sanidad.

